\begin{afterword}
\summaryhead

The post starts from a fantasy story in which a match brought from Earth will not strike in a world of Norse magic, and says: no. It then tells how Lavoisier found out what fire is. By weighing everything in a reaction, air included, Lavoisier showed that matter is neither created nor destroyed. Burning, it turned out, uses up ``vital air'' (oxygen) and makes ``fixed air'' (carbon dioxide), and that people do the same when they breathe: people run on combustion.

Matches light because of phosphorus, and phosphorus is also central to ATP, the body's energy carrier. So if a match stops working, so do you. The rules ``Matches catch fire when struck'' and ``Humans need air to breathe'' seem unconnected, and we can imagine one without the other, but that is a fact about how our minds store rules, not about the laws of nature. Phosphorus itself behaves as it does because of the laws of electrons and quarks. Reality is joined together more tightly than we like to believe.

\responsehead

The post's main point is sound. Matches and breathing do depend on shared chemistry, and a change in chemistry that stopped one would stop the other. The general lesson, that rules we store separately may be joined in nature, is the map-and-territory distinction applied to beliefs we keep apart, and the example makes it vivid.

The history is told through Lavoisier alone, although the findings had older roots. Mayow reported in 1674 that antimony heated with a burning glass gains weight, and put the gain down to particles from the air. Conservation of mass was, in the words of Wikipedia's article, ``an important assumption'' in eighteenth-century chemistry, and Lomonosov stated it in 1756. The link between burning and breathing, which the post presents as the surprise, ``had long been recognized'': Mayow had concluded that a candle and an animal use up the same part of the air, and Black that animals breathe out fixed air. Lavoisier's contribution was the chemistry and the measurements that turned an old analogy into a demonstrated process. The post's own later sentence, ``It took centuries to discover the connection,'' fits the longer history better. A smaller slip: DNA was isolated in 1869; 1953 is the year of its structure.

The link from matches to ATP depends on a different property of phosphorus than the one the post names. The post says phosphorus is highly reactive and ``thus'' suited to ATP. Elemental phosphorus reacts readily with oxygen; the phosphate in ATP is the fully oxidized form, and chemists explain life's use of it by its stability in water, not by reactivity. The post's conclusion still holds, since both matches and ATP still depend on the chemistry of one element.

In short: a sound point about the unity of natural law, told through a history that names Lavoisier alone for findings with older roots.
\end{afterword}
